\begin{table}[tb]
\centering\footnotesize\setlength{\tabcolsep}{2pt}
\caption{全24条件の成否と論理入出力token．不明生成を０へ置換せず，未評価・未実行も残す．}
\begin{tabular}{lrrr}
課題・反復 & M & H & L \\ \hline
写真移動1 & 成功 128,796 & 成功 135,628 & 成功 135,628 \\
写真移動2 & 未達 152,155 & 成功 135,628 & 成功 135,628 \\
精算1 & 成功 241,963 & 成功 211,110 & 成功 211,110 \\
精算2 & 成功 275,944 & 成功 211,110 & 成功 211,110 \\
返金1 & 成功 138,701 & 成功 102,087 & 成功 102,087 \\
返金2 & 成功 109,383 & 成功 102,087 & 成功 102,087 \\
曲情報の書出し1 & 未達 257,266 & \shortstack[r]{未達 358,638\\＋不明1} & \shortstack[r]{未達 348,419\\＋不明1} \\
曲情報の書出し2 & 未達 90,721 & \shortstack[r]{未達 358,638\\＋不明1} & \shortstack[r]{未達 346,789\\＋不明1} \\
\hline\end{tabular}
\end{table}
採用基準を満たした比較は0/4であり，両対照を両反復で上回るという固定基準は未達だった．条件を選び直した改善や，新手法の成立として扱わない．
反復1のM/H/Lは，成功／評価数が順に3/4，3/4，3/4，論理費用は766,726，807,463＋不明1，797,244＋不明1 tokenだった．
反復2のM/H/Lは，成功／評価数が順に2/4，3/4，3/4，論理費用は628,203，807,463＋不明1，795,614＋不明1 tokenだった．
\begin{table}[tb]
\centering\footnotesize\setlength{\tabcolsep}{3pt}
\caption{全４比較の費用差（L−対照）．両側成功・両側未達の寄与を分ける．不明使用量があれば確定差を求めない．}
\begin{tabular}{llrrr}
反復 & 対照 & 全課題 & 両側成功 & 両側未達 \\ \hline
1 & M & 未確定 & 未確定 & 未確定 \\
1 & H & 未確定 & 未確定 & 未確定 \\
2 & M & 未確定 & 未確定 & 未確定 \\
2 & H & 未確定 & 未確定 & 未確定 \\
\hline\end{tabular}
\end{table}
双方が未達の課題を安く終えた差は，完成に必要な費用の削減とは異なる．一方だけが成功した課題の差は全課題へ含むが，両側成功・両側未達の欄には含めない．
全予定24条件のうち，費用記録があるのは24条件，評価済みは24条件だった．新規生成は426回，既知論理入出力4,602,713 token，不明使用量4生成，実API送信は1156回だった．推論器の全稼働は16,263.578秒，うちロードは10.061秒，各仕事の初期化・作業・後片付けの合計は16,192.106秒，評価合計は41.470秒だった．これらの時間は重複し，加算しない．
